\section{Complexity of the Satisfiability Problem for $\Hxpath$}
\label{sec:complexity}

In this section, we prove that the satisfiability problem for $\Hxpath$ is \pspace-complete.
%    We obtain \pspace-hardness directly from the fact that Basic Modal Logic---a proper fragment of $\Hxpath$---is already known to be \pspace-hard (see, e.g.,~\cite{blackburn2001modal,blackburn06}).
\pspace-hardness follows from \pspace-hardness of the Basic Modal Logic~\cite{blackburn2001modal,blackburn06}, a proper fragment of $\Hxpath$.
For completeness, we present an algorithm that applies the tableau rules in a controlled manner, ensuring that satisfiability checking can be performed in polynomial space.
Our algorithm explores the model induced by the tableau in a depth-first fashion.
This requires particular care, as this model is not necessarily tree-shaped and may contain cycles or shared substructures.
% We show how to handle these complexities while remaining within \pspace, using strategies that prevent redundant revisitations and manage space usage efficiently.
%To the best of our knowledge, no detailed algorithm in the literature addresses this kind of exploration in the presence of hybrid features and data constraints. In this respect, our contribution also helps fill this gap by making 
Furthermore, we make the necessary mechanisms explicit, making our algorithm suitable for implementation, a contribution by its own sake.
% and proving that the resulting procedure runs in polynomial space.

\begin{algorithm}[t]
    \caption{$\SATALG$}\label[algorithm]{alg:start}
    \begin{small}
        \Function{$\SATALG(\varphi : \text{node expression})$}
        {
            {\bf require:} true \\
            $\NOM \gets \text{a new nominal}$\;
            $\Phi\gets [\NOM{:}\varphi]$\;
            \Return{$\SATALG(\Phi,\NOM+1)$}\;
            {\bf ensure:} $\SATALG(\varphi)$ iff $\varphi$ is satisfiable \\
        }
    \end{small}
\end{algorithm}


\begin{algorithm}[t]
    \caption{$\SATALG$}\label[algorithm]{alg:sat}
    \begin{small}
        \Function{$\SATALG(\Phi : \text{list of node expressions}, \NOM : \text{nominal})$}
        {
            {\bf require:} the nominals in $\Phi$ are in $[0 \dots \NOM)$ \\

            % {\bf var} $\STACK$, $\BRANCH$, r\;

            \If{$\Phi = [~]$}{\KwRet{true}}

            \vspace*{2pt}

            $\BRANCH, \STACK \gets \Phi, [~]$\;
            \Repeat{$\Sigma = [~]$}
            {
                % \If(\Comment*[f]{explore an alternative branch}){$\Sigma \neq [~]$}
                % {
                %     {t, $\psi$ $\gets$ $\TOP(\STACK)$}\;
                %     $\BRANCH, \STACK \gets \BRANCH[{:}\text{t}] + [\psi], \POP(\STACK)$\!
                % }
                \If(\Comment*[f]{explore an alternative branch}){$\Sigma = [(\mathrm{b},\chi)] + \mathrm{L}$}
                {
                    $\BRANCH, \STACK \gets \BRANCH_{\text{b}} + [\chi], \mathrm{L}$;
                }

                \vspace*{2pt}

                \Expand(\Comment*[f]{close $\BRANCH$ under the application of rules of type 1})
                {
                    \If{there is a rule $\RULE$ of type 1 that is untreated in $\BRANCH$}
                    {
                        {\bf let} $\Psi$ be the consequent of $\RULE$\\
                        
                        $\BRANCH \gets \BRANCH + \Psi$;
                        \Comment*[f]{expand the branch}\\
                        \GoTo expand
                    }

                    \If{there is a rule $\RULE$ of type 2 that is untreated in $\BRANCH$}
                    {
                        {\bf let} $\{\psi, \chi\}$ be the consequent of $\RULE$\\
                        
                        $\BRANCH \gets \BRANCH + [\psi]$;
                        \Comment*[f]{expand the branch}\\
                        $\STACK \gets [(|\BRANCH|, \chi)] + \Sigma$;
                        \Comment*[f]{store the alternative and its location in the branch}\\
                        \GoTo expand
                    }
                }

                \vspace*{2pt}

                \If(\Comment*[f]{try an alternative branch}){there is a clash in $\BRANCH$ and $\Sigma \neq [~]$}
                    {{\bf continue}\;}
                \ElseIf(\Comment*[f]{no alternative branches to try}){there is a clash in $\BRANCH$}
                    {\Return{false}\;}
                % \ElseIf(\Comment*[f]{$\BRANCH$ is open \& saturated}){no rule of type 3 applies in $\BRANCH$}{\KwRet{true}}

                \vspace*{2pt}

                \RecStep(\Comment*[f]{explore diamonds in a DFS manner})
                {
                    r $\gets$ true\;
                    % $\text{C}$ $\gets$ $\BRANCH'$ $+$ list of all $\lnot\tup{\alpha \cmpr \beta}$ in $\BRANCH$\;
                    \ForAll{rules $\RULE$ of type 3 applicable in $\BRANCH$}
                    {
                        {\bf let} $\psi$ be the premiss of $\RULE$ and $\psi \in \{i{:}\tup{\pathh}\chi, \tup{@_i\pathh\alpha \cmpr \beta}\}$\;

                        % {$\NOM \gets \NOM+1$};\Comment*[f]{create a new nominal}\\
                        $\Phi$ $\gets$ $\Phi$ + $\BRANCH$ minus all diamonds in $\BRANCH$ + $[i{:}\tup{\pathh}\NOM]$\;

                        % {$\Phi \gets \Phi + \UNBOX(i,\NOM,\pathh,\BRANCH'+\text{C})$};\Comment*[f]{propagate boxes}
                        \If{$\psi = i{:}\tup{\pathh}\chi$}%
                            {$\Phi \gets \Phi + [\NOM{:}\chi]$}

                        \If{$\psi = \tup{@_i\pathh\alpha \cmpr \beta}$}%
                            {$\Phi \gets \Phi + [\tup{@_{\NOM}\alpha \cmpr \beta}]$}

                        % $\Phi \gets \Phi + \UNBOX(\DIAMONDS+[i{:}\tup{\astep}\NOM],\BRANCH+\BOXES)$\;
                        {r $\gets$ r and $\SATALG(\Phi, \NOM+1)$}\;

                        $\Phi$ $\gets$ list of root literals and rooted boxes in $\Phi$\;
                    }
                }

                \vspace*{2pt}

                \lIf(\Comment*[f]{the branch is open and saturated}){r}{\Return{true}}
            }
            \KwRet{false}\;
            {\bf ensure:} $\SATALG([\NOM{:}\varphi],\NOM+1)$ iff $\varphi$ is satisfiable
        }
    \end{small}
\end{algorithm}


We begin by introducing some preliminary concepts. 
First, recall that without loss of generality, we assume that $\noms$ is the set of natural numbers. 
%We make use of (finite) lists.
In what follows we use
$[~]$ for the empty list,
$[\mathrm{e}]$ for a list with a single element $\mathrm{e}$,
$\mathrm{L}_1 + \mathrm{L}_2$ for the concatenation of lists $\mathrm{L}_1$ and $\mathrm{L}_2$,
and
$|\mathrm{L}|$ for the length %, i.e., number of elements,
of a list $\mathrm{L}$.
%All these operations return new lists.
%Regarding lists also, 
Moreover, we use $\mathrm{L} = [\mathrm{e}]+\mathrm{L}'$ to indicate a non-empty list $\mathrm{L}$ with first element $\mathrm{e}$ and tail $\mathrm{L}'$, and $\mathrm{L}_{\mathrm{n}}$ to indicate $\mathrm{L}$'s initial segment of length $\mathrm{n}$.
Finally, we assume that parameter variables are pass-by-reference and assignments are copy constructors.
% This implies that the caller and the callee use the same variable for the parameter.
% Thus, if the callee modifies the parameter variable, the effect is visible to the caller.
% Finally, we assume assignments are copy constructors.

The algorithmic application of tableau rules starts by calling the function $\SATALG$ in \Cref{alg:start} with a node expression $\varphi$.
In this function, we set up the context for the application of tableau rules.
The actual application of rules is implemented in the overloaded $\SATALG$ in \Cref{alg:sat}. %with the help of~\Cref{alg:filter}.

To keep the space complexity polynomial in \Cref{alg:sat}, we explore branches one at a time.
The information needed to explore branches is kept in local lists $\BRANCH$ and $\STACK$.
Intuitively, $\BRANCH$ stores the current branch of the tableau.
In turn, $\STACK$ behaves as a stack storing pairs $(\mathrm{b},\chi)$ where $\mathrm{b}$ is a number indicating the position to backtrack in $\BRANCH$ to explore an alternative branch, and $\chi$ is the alternative node expression to be explored.

Turning to branch expansions, we classify tableau rules into types 1, 2, and 3.
%This classification obeys the following rationale.
% \begin{description}
%     \item[(type 1)]
    The rules of type 1 are: the rules $({\lnot}{\lnot})$, $(\land)$, and $({\lnot}\Diamond)$ in \Cref{rules:basic};
    all the rules in \Cref{rules:nominals,rules:internalization};
    %and
    the rules (test), $(@)$, ($\lnot$child), $(\lnot@)$, $(\lnot{\cup})$, and (com$_1$), (com$_2$), and (dTrans) in \Cref{rules:path}; and 
    the rules
    \begin{align*}
        {\infer[(\mathrm{copy}_0)]
            {\phantom{|} k{:}\tup{\child}j}
            {\deduce[]
                {\phantom{|}i{:}j}
                {\phantom{|}k{:}\tup{\child}i}
            }
        }
        \qquad
        {\infer[(\mathrm{copy}_1)]
            {\phantom{|} \tup{@_j\alpha \cmpr \beta}}
            {\deduce[]
                {\phantom{|}i{:}j}
                {\phantom{|}\tup{@_i\alpha \cmpr \beta}}
            }
        }
        \qquad
        {\infer[(\mathrm{copy}_2)]
            {\phantom{|} \lnot\tup{@_j\alpha \cmpr \beta}}
            {\deduce[]
                {\phantom{|}i{:}j}
                {\phantom{|}\lnot\tup{@_i\alpha \cmpr \beta}}
            }
        }
    \end{align*}
    These rules have side conditions $j < i$ and $j$ a root nominal.
    In brief, these rules may be understood in analogy with (copy) in \Cref{rules:nominals}.
    These rules are admissible and do not affect termination.
    They are introduced here only to cope with space complexity.
    %
    % \item[(type 2)]
    The rules of type 2 are: the rule $(\lnot{\land})$ in \Cref{rules:basic}, and the rules ($\lnot$test), and $({\cup})$ in \Cref{rules:path}.
    %
    % \item[(type 3)]
    Finally, the rules of type 3 are: the rule $(\Diamond)$ in \Cref{rules:basic}, and
        the rule (child) in \Cref{rules:path}.
% \end{description}
The rationale behind the classification of tableau rules is that rules of type 1 extend the branch linearly, rules of type 2 split the branch into two, and rules of type 3 handle diamonds.

We are now in a position to explain how~\Cref{alg:sat} applies tableau rules to the current branch, i.e., to the list $\BRANCH$.
It is clear that the application of a rule should expand $\BRANCH$ by adding node expressions at its end.
This expansion is actually dealt with in two blocks of code:
\begin{description}
    \item[(expand)]
        Lines 9 to 18.
        This block of code applies rules of type 1 or of type 2 until there are no untreated rules.
        The case of a rule $\RULE$ of type 1 is straightforward.
        Let $\Psi$ be the set of node expressions in the consequent of $\RULE$.
        We say that $\RULE$ is untreated iff $\Psi \nsubseteq \BRANCH$.
        If $\RULE$ is untreated, we add all the node expressions in $\Psi$ at the end of $\BRANCH$. 
        %In the case of a rule $\RULE$ of type 2.
        In turn, let $\RULE$ be a rule of type 2 and let $\Psi = \{\psi,\chi\}$ be the node expressions in the consequent of $\RULE$. 
        We say that $\RULE$ is untreated iff $\Psi \cap \BRANCH = \emptyset$.
        If $\RULE$ is untreated, we add $\psi$ at the end of $\BRANCH$, and store in $\STACK$ the alternative $\chi$ together with the point in $\BRANCH$ at which $\psi$ is introduced.
        % The information in $\STACK$ enables us to later on ``backtrack'' and explore an alternative branch.
        % Upon termination of this loop, we can check if there is a clash in $\BRANCH$.
        % If there is a clash and there is an alternative branch, then, we backtrack and explore an alternative branch.
        % The check is implemented in line 17 and the backtracking, which uses the information stored in $\STACK$, is implemented in lines 7 to 8.
        % If there is a clash and there is no alternative branch to try, then, the node expression is unsatisfiable

    \item[(explore)]
        Lines 23 to 33.
        This block deals with the exploration of ``diamonds'', i.e., node expressions of the form $i{:}\tup{\pathh}\psi$ that are not accesibility constraints, or node expressions of the form $\tup{@_i\pathh\alpha \cmpr \beta}$. 
        To do so, we take some ideas from \cite{Ladner77,Marx06} of exploring diamonds in a depth first search (DFS) manner and adapt them to our setting.
        
        First, in our case, the exploration of diamonds corresponds to the application of rules of type 3.
        A rule $\RULE$ of type 3 is applicable iff its premiss belongs to $\BRANCH$.
        If $\RULE$ is applicable, its application involves: creating a new nominal $\NOM$, creating the accesibility constraint $i{:}\tup{\child}\NOM$, and adding the node expression in the consequent of $\RULE$ to the branch.

        Second, exploring diamonds in a DFS manner involves recursive calls to $\SATALG$.
        These recursive calls take into account all alternative explorations of diamonds and put a temporary pause on them. Notice that only one diamond must be treated at each step, as they can be individually satisfied and only can find contradictions with box expressions. Once the satisfiability of one particular diamond is checked, the algorithm backtracks and all diamonds are restored on $\BRANCH$. 
        This is implemented by removing all alternative diamonds from the branch in line 27 of \Cref{alg:sat}, before the recursive call.

        Lastly, before backtracking in the DFS strategy we may need to gather \emph{global} information to be passed along to alternative calls to $\SATALG$.
        This global information %is \emph{global} to the branch of the tableau and 
        is mediated by root nominals in a tableau branch and corresponds to \emph{root literals} and \emph{rooted boxes}.
        By a literal, we mean a node expression of the form $i{:}j$, $i{:}\lnot j$, $i{:}p$, $i{:}\lnot p$, $i{:}\tup{\child}j$, $\tup{@_i \cmpr @_j}$ where $\{i,j\} \subseteq \noms$ and $p \in \props$.
        A root literal is one whose nominals are root nominals.
        Root literals contain the relevant information stored in urfathers of root nominals, thus they need to be passed to the next recursive call.
        A rooted box is a node expression of the form $i{:}\lnot\tup{\child}\psi$ or $\lnot\tup{@_i\alpha \cmpr @_j\beta}$ where $i$ and $j$ are root nominals.
        Global information to the branch is recorded in line 33 of \Cref{alg:sat}.
\end{description}


Correctness of \Cref{alg:sat} follows from the correctness of tableau rules, while termination is ensured since each loop and recursive call is bounded. The next result shows the algorithm uses only polynomial space.

\begin{thm}
    \label{th:satpspace}
    The satisfiability problem for $\Hxpath$ is \pspace-complete.
\end{thm}
\begin{proof}
    Hardness follows from the fact that the Basic Modal Logic~$\K$ is \pspace-hard~\cite{blackburn2001modal,blackburn06}, and $\Hxpath$ subsumes it. 
    Completeness follows from the fact that \Cref{alg:sat} uses at most polynomial space with respect to the size of the input node expression, as analyzed in what follows. 

    The local variables $\BRANCH$ and $\STACK$ occupy at most a polynomial amount of memory (bounded by the number of untreated rules of type 1 and 2).
    The number of recursive calls, in a branch, is also polynomial (bounded by the number of nested diamonds, including path expressions of the form $\child$). 
    The only information that is global to the branch is that on the list $\Phi$.
    But this list is also of polynomial size, as the only persistent information that is passed along between different recursive calls refers only to root literals, and this information is polynomial.

    When put together, these facts tell us that \Cref{alg:sat} uses at most a polynomial amount of memory with respect to the size of the input node expression.
\end{proof}

We close this section with a run of the function $\SATALG$.
This run illustrates the inner workings of the algorithm it implements.
Moreover, it allows us to recognize how it uses only polynomial space.

\begin{exa}
Let us suppose that we start $\SATALG$ in~\Cref{alg:start} with the node expression
\[
    \varphi = \tup{\child(1\land p)? =_\cmp \astep[b]} \land \tup{\astep[c]@_1(\neg p)? =_\cmp \astep[b]}  \land \neg\tup{\astep[b]\neq_\cmp @_1}.
\]   
The call to $\SATALG(\varphi)$ invokes $\SATALG(\Phi,\NOM)$ with $\Phi=[2{:}\varphi]$ and  $\NOM=3$.
We proceed to list the different steps in~\Cref{alg:sat} on this input.
Sometimes we group several steps together and omit some node expressions to make the presentation more succinct. 

\begin{description}
  \item[Initialization] 
    We set $\BRANCH=\Phi$ and $\STACK=[~]$.
    Notice that the sole root nominal in $\varphi$ is $1$.

  \item[Step 1] 
    We reach the `expand' block and apply rules of type 1 until no such rule is left untreated.
    The list $\BRANCH$ is updated during this process (the rules applied to obtain the node expressions are indicated on the right):
    \[
      \begin{array}{r@{\,}@{}lr}
        \BRANCH =
          & [
              2{:}\tup{\child(1\land p)? =_\cmp \astep[b]},
              2{:}\tup{\astep[c]@_1(\neg p)? =_\cmp \astep[b]},
              2{:}\neg\tup{\astep[b]\neq_\cmp @_1},
          ] + \dots +
              & (\land) \\
          & [
              \tup{@_2\child(1\land p)? =_\cmp @_2\astep[b]},
              \tup{@_2\astep[c]@_1(\neg p)? =_\cmp @_2\astep[b]},
              \neg\tup{@_2\astep[b]\neq_\cmp @_2@_1}
            ]
              & (\text{int}_1), (\text{int}_2)
      \end{array}
    \]

  \item[Step 2]
    We reach the `explore' block and choose to apply the rule of type 3 with premiss $\tup{@_2\child(1\land p)? =_\cmp @_2\astep[b]}$.
    Before the recursive call we get
    \[
      \Phi =
        [
          \neg\tup{@_2\astep[b]\neq_\cmp @_2@_1}, 2{:}\tup{\astep[a]}3, \tup{@_3(1\land p)? =_\cmp @_2\astep[b]}
        ]
    \]

  \item[Step 3] 
    We call $\SATALG(\Phi,\NOM)$ recursively for a first time, with $\NOM=4$. 
    In the `expand' block of the first recursive call, we apply all rules of type 1 and get
    \[
      \begin{array}{r@{\,}@{}lr}
        \BRANCH = 
          &
          [
            \neg\tup{@_2\astep[b]\neq_\cmp @_2@_1},
            2{:}\tup{\astep[a]}3,
            \tup{@_3(1\land p)? =_\cmp @_2\astep[b]},
          ] + \dots + \\
          &
          [
            3{:}1,3{:}p,1{:}3,1{:}p, 2{:}\tup{\astep[a]}1,
            \tup{@_3 =_\cmp @_2\astep[b]},
            \tup{@_1 =_\cmp @_2\astep[b]},
            \tup{@_2\astep[b]=_\cmp@_1}
          ]
      \end{array}
    \]
    Notice that $1{:}p$ is obtained from (copy), $2{:}\tup{\child}1$ from (copy$_0$) and $\tup{@_1=_\cmp @_2\astep[b]}$ from (copy$_1$). 
    This way, the nominal $3$ is replaced by the root nominal $1$.
    Finally, the last node expression in $\BRANCH$ is obtained by applying the rule (com$_1$).

  \item[Step 4]
    In the `explore' block of the first recursive call we choose to apply the rule of type 3 with premiss $\tup{@_2\astep[b]=_\cmp@_1}$.
    Before the recursive call we get
    \[
      \Phi =
        [
          \neg\tup{@_2\astep[b]\neq_\cmp @_2@_1}, 1{:}p, 2{:}\tup{\astep[a]}1, 2{:}\tup{\astep[b]}4, \tup{@_4 =_{\cmp} @_1}
        ]
    \]

  \item[Step 5]
    We call $\SATALG(\Phi,\NOM)$ recursively for a second time (with $\NOM=5$).
    In the `expand' block of the second recursive call, we apply all rules of type 1 and get
    \[
      \begin{array}{r@{\,}@{}lr}
        \BRANCH = 
            &
            [
              \neg\tup{@_2\astep[b]\neq_\cmp @_2@_1},
              2{:}\tup{\astep[b]}4,
              \tup{@_4 =_{\cmp} @_2@_1}
            ] + \dots + \\
            &
            [
              \neg \tup{@_4 \neq_{\cmp} @_2@_1}
            ] + \dots +
              & \text{($\lnot$child)} \\
            &
            [
              \neg \tup{@_4 \neq_{\cmp} @_1}
            ]
              & \text{($\neg@$)}
      \end{array}
    \]
        

  \item[Step 6] 
    Since no rules of type 2 are triggered, no clash is found, and there are some untreated node expressions, the algorithm proceeds to the `explore' phase again. A new rule of type 3 is triggered, this time with premise $\psi= \tup{@_2\astep[c]@_1(\neg p)? =_\cmp @_2\astep[b]}$. Before the new recursive call, we get
    \[
      \Phi =
        [
          \neg\tup{@_2\neq_{\cmp}@_1}, 1{:}p
        ] 
        + \dots + 
        [
          2{:}\tup{\astep[c]}5, \tup{@_5@_1(\neg p)? =_\cmp @_2\astep[b]} 
        ]
    \]


  \item[Step 7] 
    The recursive call invokes $\SATALG(\Phi,\NOM)$, with $\NOM=6$. In the `expand' block, we get 
    \[
      \begin{array}{r@{\,}@{}lr}
        \BRANCH = 
            &
            [
              \neg\tup{@_2\neq_{\cmp}@_1}, 1{:}p, 2{:}\tup{\astep[c]}5, \tup{@_5@_1(\neg p)? =_\cmp @_2\astep[b]}
            ] + \dots + \\
            &
            [
              \tup{@_1(\neg p)? =_{\cmp} @_2\astep[b]}, \tup{@_1 =_{\cmp}@_2\astep[b]}, 1{:}\neg p
            ] 
              & (@), (?) 
      \end{array}
    \]
    \item[Step 8] 
      After executing the Step 7, the algorithm finds a clash in line 21, since both $1{:}p$ and $1{:}\neg p$ belong to the sole tableau branch. Therefore, \Cref{alg:sat} returns that $\varphi$ is unsatisfiable.
\end{description}

Notice how in each step the algorithm keeps track of the list of accesibility constraints created in the exploration of diamonds in a DFS manner.
The length of these lists is bounded by the lengths of the path expressions traversed in each diamond.
Moreover, we only need to keep track of a polynomially bounded number of node expressions at each point in these paths.
This example illustrates the polynomial space bound in \Cref{th:satpspace}.
\end{exa}